% GENERATED FILE. Edit canonical lessons, then run npm run generate:books.
% canonical-source-hash: fnv1a64:386e8ca4b566cba1
% canonical-lessons: TA-C543-kattalai, TA-C543-katumaiyana, TA-C543-kairekai, TA-C543-taniyurimai, TA-C543-coti

\chapter{Rules and Checks}
\label{ch:ta-rules-and-checks}
\hlchaptermodality{543}

\begin{chapteropening}
\textbf{By the end of this chapter:} \emph{I can talk about orders, strict rules and checks.}

Complete the last lesson of chapter 543: I can talk about orders, strict rules and checks.
\end{chapteropening}

\section[\texorpdfstring{\ta{கட்டளை} (kaṭṭaḷai)}{கட்டளை (kaṭṭaḷai)}]{\ta{கட்டளை} (kaṭṭaḷai) — an official order; a command}

\label{lesson:TA-C543-kattalai}



\begin{quote}
Before the new one: say the Tamil for to swear an oath, then the Tamil for to grant.
\end{quote}

\subsection*{You'll want to know: \ta{கட்டளை}}
\textbf{\ta{கட்டளை}} — \emph{kaṭṭaḷai} — \textquotedblleft{}an official order; a command\textquotedblright{}.

\begin{quote}
\ta{அதிகாரியின்} \ta{கட்டளையை} \ta{எல்லோரும்} \ta{பின்பற்றினார்கள்}.
\end{quote}

\begin{quote}
\emph{atikāriyiṉ kaṭṭaḷaiyai ellōrum piṉpaṟṟiṉārkaḷ.}
\end{quote}

\begin{quote}
Everyone followed the officer's order.
\end{quote}

In computing it also means a command. To give an order is \ta{கட்டளை} \ta{இடு}.

\begin{cousinweb}
\textbf{\ta{கட்டளை}} (\emph{kaṭṭaḷai}) is native Dravidian: DEDR 1147, under \ta{கட்டு} (\emph{kaṭṭu}), \textquotedblleft{}to tie, to bind\textquotedblright{}, gives \ta{கட்டளை} (\emph{kaṭṭaḷai}), \textquotedblleft{}a code, a rule, regulations\textquotedblright{}.
\end{cousinweb}

\begin{grammarlens}[title={What you've built}]
One more word for saying what happens if.
\end{grammarlens}

\subsection*{Your turn}
\begin{itemize}
\raggedright
  \item \emph{Say it:} \emph{kaṭṭaḷai}
  \item \emph{Say it:} \emph{kaṭṭaḷai}, once more
  \item \emph{Say it:} \emph{atikāriyiṉ kaṭṭaḷaiyai ellōrum piṉpaṟṟiṉārkaḷ.} — Everyone followed the officer's order.
  \item \emph{Recall:} say the Tamil for to swear an oath, then the Tamil for to grant, then say \emph{kaṭṭaḷai} again
\end{itemize}

\begin{tcolorbox}[breakable,colback=teal!4,colframe=teal!35!black,title={Before you move on}]
What does \ta{கட்டளை} mean? (\textquotedblleft{}An official order; a command\textquotedblright{}.) Say it once more.
\end{tcolorbox}
\section[\texorpdfstring{\ta{கடுமையான} (kaṭumaiyāṉa)}{கடுமையான (kaṭumaiyāṉa)}]{\ta{கடுமையான} (kaṭumaiyāṉa) — strict; severe}

\label{lesson:TA-C543-katumaiyana}



\begin{quote}
Before the new one: say the Tamil for to grant, then the Tamil for an official order.
\end{quote}

\subsection*{You'll want to know: \ta{கடுமையான}}
\textbf{\ta{கடுமையான}} — \emph{kaṭumaiyāṉa} — \textquotedblleft{}strict; severe\textquotedblright{}.

\begin{quote}
\ta{விதிகளை} \ta{மீறினால்} \ta{கடுமையான} \ta{அபராதம்} \ta{உண்டு}.
\end{quote}

\begin{quote}
\emph{vitikaḷai mīṟiṉāl kaṭumaiyāṉa aparātam uṇṭu.}
\end{quote}

\begin{quote}
If you break the rules, there is a severe fine.
\end{quote}

From the noun \ta{கடுமை}, harshness, with -\ta{ஆன}. As an adverb: \ta{கடுமையாக} \ta{உழை}, to work hard.

\begin{cousinweb}
\textbf{\ta{கடுமையான}} (\emph{kaṭumaiyāṉa}) is native Dravidian, the -\ta{ஆன} adjective of \ta{கடுமை} (\emph{kaṭumai}), \textquotedblleft{}severity, harshness\textquotedblright{}, which DEDR 1135 gives under \ta{கடு}, \textquotedblleft{}to throb, to smart\textquotedblright{}.
\end{cousinweb}

\begin{grammarlens}[title={What you've built}]
One more word for saying what happens if.
\end{grammarlens}

\subsection*{Your turn}
\begin{itemize}
\raggedright
  \item \emph{Say it:} \emph{kaṭumaiyāṉa}
  \item \emph{Say it:} \emph{kaṭumaiyāṉa}, once more
  \item \emph{Say it:} \emph{vitikaḷai mīṟiṉāl kaṭumaiyāṉa aparātam uṇṭu.} — If you break the rules, there is a severe fine.
  \item \emph{Recall:} say the Tamil for to grant, then the Tamil for an official order, then say \emph{kaṭumaiyāṉa} again
\end{itemize}

\begin{tcolorbox}[breakable,colback=teal!4,colframe=teal!35!black,title={Before you move on}]
What does \ta{கடுமையான} mean? (\textquotedblleft{}Strict; severe\textquotedblright{}.) Say it once more.
\end{tcolorbox}
\section[\texorpdfstring{\ta{கைரேகை} (kairēkai)}{கைரேகை (kairēkai)}]{\ta{கைரேகை} (kairēkai) — a fingerprint}

\label{lesson:TA-C543-kairekai}



\begin{quote}
Before the new one: say the Tamil for an official order, then the Tamil for strict.
\end{quote}

\subsection*{You'll want to know: \ta{கைரேகை}}
\textbf{\ta{கைரேகை}} — \emph{kairēkai} — \textquotedblleft{}a fingerprint\textquotedblright{}.

\begin{quote}
\ta{காவலர்கள்} \ta{கதவில்} \ta{இருந்த} \ta{கைரேகையைப்} \ta{பதிவு} \ta{செய்தார்கள்}.
\end{quote}

\begin{quote}
\emph{kāvalarkaḷ katavil irunta kairēkaiyaip pativu ceytārkaḷ.}
\end{quote}

\begin{quote}
The police recorded the fingerprint on the door.
\end{quote}

Also the lines of the palm. \ta{ரேகை} alone means \textquotedblleft{}a line\textquotedblright{}.

\begin{cousinweb}
\textbf{\ta{கைரேகை}} (\emph{kairēkai}) joins the native \ta{கை} (\emph{kai}), \textquotedblleft{}a hand\textquotedblright{}, and \ta{ரேகை} (\emph{rēkai}), \textquotedblleft{}a line\textquotedblright{}, from Sanskrit \emph{rekhā}, \textquotedblleft{}a scratch, a line\textquotedblright{} (CDIAL 10810).
\end{cousinweb}

\begin{grammarlens}[title={What you've built}]
One more word for saying what happens if.
\end{grammarlens}

\subsection*{Your turn}
\begin{itemize}
\raggedright
  \item \emph{Say it:} \emph{kairēkai}
  \item \emph{Say it:} \emph{kairēkai}, once more
  \item \emph{Say it:} \emph{kāvalarkaḷ katavil irunta kairēkaiyaip pativu ceytārkaḷ.} — The police recorded the fingerprint on the door.
  \item \emph{Recall:} say the Tamil for an official order, then the Tamil for strict, then say \emph{kairēkai} again
\end{itemize}

\begin{tcolorbox}[breakable,colback=teal!4,colframe=teal!35!black,title={Before you move on}]
What does \ta{கைரேகை} mean? (\textquotedblleft{}A fingerprint\textquotedblright{}.) Say it once more.
\end{tcolorbox}
\section[\texorpdfstring{\ta{தனியுரிமை} (taṉiyurimai)}{தனியுரிமை (taṉiyurimai)}]{\ta{தனியுரிமை} (taṉiyurimai) — privacy}

\label{lesson:TA-C543-taniyurimai}



\begin{quote}
Before the new one: say the Tamil for strict, then the Tamil for a fingerprint.
\end{quote}

\subsection*{You'll want to know: \ta{தனியுரிமை}}
\textbf{\ta{தனியுரிமை}} — \emph{taṉiyurimai} — \textquotedblleft{}privacy\textquotedblright{}.

\begin{quote}
\ta{இந்தச்} \ta{செயலி} \ta{உங்கள்} \ta{தனியுரிமையை} \ta{மதிக்கிறது}.
\end{quote}

\begin{quote}
\emph{intac ceyali uṅkaḷ taṉiyurimaiyai matikkiṟatu.}
\end{quote}

\begin{quote}
This app respects your privacy.
\end{quote}

\ta{தனி} \textquotedblleft{}alone, private\textquotedblright{} + \ta{உரிமை} (\emph{urimai}) \textquotedblleft{}a right\textquotedblright{}.

\begin{cousinweb}
\textbf{\ta{தனியுரிமை}} (\emph{taṉiyurimai}) is a modern native compound: \ta{தனி} (\emph{taṉi}), \textquotedblleft{}alone, single, private\textquotedblright{}, and \ta{உரிமை} (\emph{urimai}), \textquotedblleft{}a right, ownership\textquotedblright{} (DEDR 654).
\end{cousinweb}

\begin{grammarlens}[title={What you've built}]
One more word for saying what happens if.
\end{grammarlens}

\subsection*{Your turn}
\begin{itemize}
\raggedright
  \item \emph{Say it:} \emph{taṉiyurimai}
  \item \emph{Say it:} \emph{taṉiyurimai}, once more
  \item \emph{Say it:} \emph{intac ceyali uṅkaḷ taṉiyurimaiyai matikkiṟatu.} — This app respects your privacy.
  \item \emph{Recall:} say the Tamil for strict, then the Tamil for a fingerprint, then say \emph{taṉiyurimai} again
\end{itemize}

\begin{tcolorbox}[breakable,colback=teal!4,colframe=teal!35!black,title={Before you move on}]
What does \ta{தனியுரிமை} mean? (\textquotedblleft{}Privacy\textquotedblright{}.) Say it once more.
\end{tcolorbox}
\section[\texorpdfstring{\ta{சோதி} (cōti)}{சோதி (cōti)}]{\ta{சோதி} (cōti) — to test; to check}

\label{lesson:TA-C543-coti}



\begin{quote}
Before the new one: say the Tamil for a fingerprint, then the Tamil for privacy.
\end{quote}

\subsection*{You'll want to know: \ta{சோதி}}
\textbf{\ta{சோதி}} — \emph{cōti} — \textquotedblleft{}to test; to check\textquotedblright{}.

\begin{quote}
\ta{பாதுகாவலர்} \ta{எல்லாப்} \ta{பைகளையும்} \ta{சோதித்தார்}.
\end{quote}

\begin{quote}
\emph{pātukāvalar ellāp paikaḷaiyum cōtittār.}
\end{quote}

\begin{quote}
The guard checked all the bags.
\end{quote}

Past (I …): \ta{சோதித்தேன்}, \emph{cōtittēṉ}.

What is checked takes -\ta{ஐ}. The noun is \ta{சோதனை} (\emph{cōtaṉai}), a test.

\begin{cousinweb}
\textbf{\ta{சோதி}} (\emph{cōti}) is borrowed from the Sanskrit stem of \emph{śōdhana}, \textquotedblleft{}cleaning, purifying\textquotedblright{} (CDIAL 12629), made a Tamil verb with -\ta{இ}.
\end{cousinweb}

\begin{grammarlens}[title={What you've built}]
One more word for saying what happens if.
\end{grammarlens}

\subsection*{Your turn}
\begin{itemize}
\raggedright
  \item \emph{Say it:} \emph{cōti}
  \item \emph{Say it:} \emph{cōti}, once more
  \item \emph{Say it:} \emph{pātukāvalar ellāp paikaḷaiyum cōtittār.} — The guard checked all the bags.
  \item \emph{Recall:} say the Tamil for a fingerprint, then the Tamil for privacy, then say \emph{cōti} again
\end{itemize}

\begin{tcolorbox}[breakable,colback=teal!4,colframe=teal!35!black,title={Before you move on}]
What does \ta{சோதி} mean? (\textquotedblleft{}To test; to check\textquotedblright{}.) Say it once more.
\end{tcolorbox}
